%! Sinusoidal Functions — Unit 6, Lesson 6.6 — Homework (Answer Key)
\documentclass[10pt]{article}
\usepackage{precalculus-article}
\usepackage{precalculus-key}
\newcommand{\TallMath}[1]{$\displaystyle #1\rule[-1.4em]{0pt}{3.2em}$}

\begin{document}

\pageheader{Unit 6, Lesson 6.6}{Homework --- Answer Key}
\namedateperiod

\vspace{0.12in}

\begin{notesbox}{Sinusoidal Functions --- Practice --- Key}

\textbf{1.} Max $= 10$, min $= -2$.

\medskip
(a) Amplitude $= \dfrac{10 - (-2)}{2} = \dfrac{12}{2} = $ \ans{6}

\medskip
(b) Midline: $y = \dfrac{10 + (-2)}{2} = \dfrac{8}{2} = $ \ans{$y = 4$}

\vspace{0.1in}

\textbf{2.} Period $= 0.002$ s.

\medskip
(a) Frequency $= \dfrac{1}{0.002} = $ \ans{500 Hz}

\medskip
(b) Period of 250 Hz wave $= \dfrac{1}{250} = $ \ans{0.004 seconds}

\vspace{0.1in}

\textbf{3.} $f(\theta) = \sin(2\theta)$.

\medskip
$f(-\theta) = \sin(2(-\theta)) = \sin(-2\theta)$.

Because sine is an odd function, $\sin(-2\theta) = -\sin(2\theta) = -f(\theta)$.

Therefore $f(-\theta) = -f(\theta)$ for all $\theta$, so $f$ is \ans{odd.}

\vspace{0.1in}

\textbf{4.} AP-Style Multiple Choice.

Midline $= \dfrac{6 + (-2)}{2} = \dfrac{4}{2} = 2$, so midline is $y = 2$.

Answer: \ans{(B)}

\vspace{0.1in}

\textbf{5.} Frequency $= \dfrac{3}{2\pi}$.

Period $= \dfrac{1}{\text{frequency}} = \dfrac{1}{3/(2\pi)} = \dfrac{2\pi}{3}$,
so period $= $ \ans{$\dfrac{2\pi}{3}$}

\vspace{0.1in}

\textbf{6.} Amplitude $= 5$, midline $y = -3$.

Maximum $= $ midline $+$ amplitude $= -3 + 5 = $ \ans{2}

Minimum $= $ midline $-$ amplitude $= -3 - 5 = $ \ans{$-8$}

\end{notesbox}

\vspace{0.08in}

\begin{extensionbox}
\textbf{Extension --- Key}

If $f(-\theta) = -f(\theta)$ and $f(-\theta) = f(\theta)$ both hold, then
$f(\theta) = -f(\theta)$, which means $2f(\theta) = 0$, so $f(\theta) = 0$ for
all $\theta$.

\ans{The function must be the zero function (output $= 0$ for every input).
Its amplitude is $0$ and its midline is $y = 0$.}
\end{extensionbox}

\vspace{0.08in}

\begin{tcolorbox}[colback=goldbg, colframe=goldacc, arc=2mm,
    left=3mm, right=3mm, top=2mm, bottom=2mm]
\textbf{Preview of Lesson 3.6 --- Sinusoidal Function Transformations}

In the next lesson you will explore the general form
$g(\theta) = a\sin(b(\theta + c)) + d$.
The parameters $a$, $b$, $c$, and $d$ each control one of the characteristics
you studied today. As you finish this homework, think about which parameter
might control amplitude, which controls period, and which determines the midline.
\end{tcolorbox}

\end{document}
